% TOP_PROBLEM: {"id": 1395, "title": "Cereceda's Conjecture", "category_id": 3, "category": {"id": 3, "name": "graph_theory", "display_name": "Graph Theory", "description": "Problems involving graphs, networks, and their properties.", "slug": "graph-theory", "order_index": 3, "created_at": "2026-07-31T15:26:25.671Z"}, "status": "open"}
% FIRST_POSTED: null
% ENTRY_KIND: statement_only
% Imported from: batch06.tex; UnsolvedMath ID 1395
\documentclass[11pt]{article}
\usepackage[margin=25mm]{geometry}
\usepackage{fontspec}
\setmainfont{Latin Modern Roman}
\usepackage{amsmath,amssymb,mathtools}
\usepackage{unicode-math}
\setmathfont{Latin Modern Math}
\usepackage{xurl,hyperref}
\hypersetup{colorlinks=true,urlcolor=blue}
\setcounter{secnumdepth}{0}
\setlength{\parindent}{0pt}
\setlength{\parskip}{5pt}
\setlength{\emergencystretch}{3em}
\newcommand{\sref}[2]{\hyperlink{source:#1:#2}{[S#2]}}
\newcommand{\cataloguescope}[1]{\paragraph{Recorded target.}#1}
\begin{document}
% UnsolvedMath ID 1395; source catalogue code GRAPH-008
\setcounter{section}{1394}
\section{Cereceda's Conjecture}\label{1395}
{\small\sffamily UnsolvedMath ID 1395 · Graph Theory}
\subsection{Definitions and mathematical statement}

\paragraph{Shared notation.}$\mathbb N=\{1,2,\ldots\}$, and $\mathbb Z,\mathbb Q,\mathbb R,\mathbb C$ denote the integers, rationals, reals and complex numbers. Any other conventions are specified locally.


A finite simple graph $G$ is $d$-degenerate if every nonempty subgraph has a vertex of degree at most $d$. A proper $q$-coloring is a map $c:V(G)\to\{1,\ldots,q\}$ whose values differ on adjacent vertices. The reconfiguration graph $\mathcal R_q(G)$ has these proper colorings as vertices; two are adjacent when they differ at exactly one vertex. Its distance counts single-vertex recolorings that preserve propriety at every step.
\paragraph{Statement recorded in the supplied source.}
For every fixed integer $d\ge0$, do there exist constants $C_d>0$ and $a_d>0$ such that every $n$-vertex $d$-degenerate graph has
\[
 \operatorname{diam}\mathcal R_{d+2}(G)\le C_d n^{a_d}?
\]
The assertion includes connectivity of the reconfiguration graph. \sref{1395}{1}\sref{1395}{2}
This polynomial formulation is established by the cited bound $O(n^{d+1})$; the stronger customary quadratic bound $O(n^2)$ is a distinct question.
\subsection{Short English statement}
For each fixed degeneracy, can any two proper colorings using two extra colors be joined by a polynomial-length sequence of valid single-vertex recolorings?
\subsection{Sources}
\begin{description}
\item[\hypertarget{source:1395:1}{S1}] ULAM.AI and the UnsolvedMath Contributors, \emph{UnsolvedMath: A Curated Collection of Open Mathematics Problems}, record 1395 (GRAPH-008). CC BY 4.0. \url{https://huggingface.co/datasets/ulamai/UnsolvedMath}.
\item[\hypertarget{source:1395:2}{S2}] Nicolas Bousquet and Marc Heinrich, \emph{A polynomial version of Cereceda's conjecture}, Journal of Combinatorial Theory, Series B 155 (2022); arXiv:1903.05619v1, abstract and Section 1. \url{https://arxiv.org/abs/1903.05619v1}.
\end{description}
\label{end:1395}
\end{document}
